\chapter{Tehnologii și Framework-uri utilizate}

\section{Framework Backend: .NET and ASP.NET Core}

\subsection{Ecosistemul C\# si .NET 8}
Alegerea folosirii limbajului C\# și a platformei .NET 8 a fost motivată de necesitatea unui ecosistem robust, capabil să susțină o logică de business complexă pentru aplicații de tip întreprindere.Aplicația impune implementarea unor algoritmi avansați pentru orchestrarea datelor, după cum urmează:

\begin{itemize}
    \item  \textbf{Optimizarea algoritmică a itinerariului:} implică rezolvarea unor variații ale problemei comis-voiajorului și aplicarea unor euristici de rutare. Sistemul calculează traiectoriile optime între locații (noduri) multiple, respectând riguros constrângerile multidimensionale de distanță fizică și ferestrele de timp.
    \item  \textbf{Mecanisme de potrivire:} presupun filtrarea avansată și ierarhizarea ponderată a punctelor de interes. Algoritmul procesează preferințele utilizatorului pe trei nivele de prioritate, realizând o mapare semantică între locații și vectorul de preferințe.
    \item  \textbf{Programarea temporală:} soluționarea unor probleme de satisfacere a constrângerilor pentru alocarea și gestionarea intervalelor zilnice, asigurând crearea unor planificări fără suprapuneri, cu timpi de tranzit calculați și cu o marjă de siguranță pentru rezervări, acestea având cea mai mare prioritate.

\end{itemize}

C\# oferă avantajul verificării statice a tipurilor la compilare (static type checking) și al un sistem de tipuri puternice (strong type). Aceste caracteristici scad semnificativ riscul excepțiilor la execuție (runtime exceptions) și garantează integritatea memoriei, un atribut critic pentru consistența funcțională a aplicației.

În plus, trecerea la .NET 8 introduce optimizări la nivel de execuție, îmbunătățiri ale compilatorului JIT (Just-in-Time) și opțiuni de compilare AOT (Ahead of Time)\href{https://medium.com/@Mr.Jafri_01/just-in-time-jit-vs-ahead-of-time-aot-compilation-in-net-core-c7650f1a366a}{}\cite{jitAOT_ref}. Acestea se traduc printr-un timp de inițializare extrem de redus și o amprentă de memorie scăzută, transformând aplicația într-o soluție ideală pentru implementări scalabile de tip cloud-native. În același context, extensia LINQ (Language Integrated Query) a fost esențială în arhitectura sistemului, permițând manipularea declarativă și evaluarea întârziată(lazy evaluation) a colecțiilor în memorie, esențiale pentru modulele de filtrare complexă de filtrare și potrivire(matching).\href{https://learn.microsoft.com/en-us/dotnet/standard/linq/deferred-execution-lazy-evaluation?utm_source=chatgpt.com}{}\cite{linq_ref}

\subsection{Arhitectura ASP.NET Core}

ASP.NET Core este un framework web multiplatformă de înaltă performanță, asincron și puternic optimizat pentru procesarea eficientă a unui volum ridicat de cereri concurente, care suportă atât tiparul MVC, cât și Razor Pages.

Pentru a respecta principiul responsabilităților unice (Single Responsibility Principle — SRP), s-a optat pentru un design hibrid care combină paradigma MVC (Model-View-Controller) cu Razor Pages:

\begin{itemize}
    \item Controller: reprezintă nivelul de interacțiune fără stare (stateless), expunând puncte terminale pentru interfețele client. Acestea încapsulează și deleagă către servicii logica de afaceri(gruparea călătoriilor, managementul rezervărilor, coordonarea locațiilor)\href{https://learn.microsoft.com/en-us/aspnet/core/mvc/controllers/actions?view=aspnetcore-10.0}{ }\cite{controller_ref}.
    \item Razor Pages: găzduiesc subsistemul de Identity and Access Management (IAM), izolându-l funcțional. Această abordare centrată pe pagină (page-focused) a oferit un mediu coerent și sigur pentru gestionarea fluxului de autentificare și autorizare, fiind un ecosistem deja maturizat și pre-testat împotriva vulnerabilităților web comune (CSRF, XSS)\href{https://learn.microsoft.com/en-us/aspnet/core/razor-pages/security/authorization/simple?view=aspnetcore-10.0&utm_source=chatgpt.com}{}\cite{razor_ref}.

\end{itemize}


\textbf{Pipeline-ul middleware (infrastructura de procesare):}

Aplicația utilizează un pipeline HTTP configurabil, bazat pe lanțul de responsabilitate (Chain of Responsibility). Aceste pipeline-uri permit interceptarea și procesarea fluxului HTTP/HTTPS prin module intermediare injectate secvențial, asigurând funcții critice precum: terminarea TLS (HTTPS redirection), negocierea resurselor statice, validarea schemelor de autentificare (ex. JWT sau Cookies) și tratarea centralizată a excepțiilor globale.

\textbf{Inversion of Control (IoC) și Dependency Injection (DI):}

De asemenea utilizează containerul nativ de IoC (prin interfața IServiceCollection) care facilitează implementarea conceptelor de arhitectură curată prin injectarea dependențelor, permițând astfel ca interfețele să fie rezolvate automat cu implementările lor concrete la momentul execuției. Serviciile de infrastructură precum IHttpClientFactory, pentru apeluri externe optimizate și evitarea epuizării socket-urilor, sau IEmailSender au fost integrate fluent folosind acest container nativ, reducând cuplarea.\href{https://learn.microsoft.com/en-us/aspnet/core/fundamentals/dependency-injection?view=aspnetcore-10.0&utm_source=chatgpt.com}{}\cite{dependency_ref}

\section{Persistența datelor}

\subsection{Strategia ORM cu Entity Framework Core}
Entity Framework Core este un instrument de mapare obiect-relațional (ORM) care abstractizează operațiile bazei de date, permițând dezvoltatorilor să lucreze cu modele puternic tipizate, eliminând necesitatea interacțiunii directe prin interogări SQL brute. În contextul acestei aplicații, EF Core oferă o serie de avantaje arhitecturale.

Gestionarea relațiilor: Arhitectura datelor din aplicație prezintă interdependențe elaborate:

\begin{itemize}
    \item  One-to-Many: Utilizatori $\to$ Călătorii $\to$ Planuri Zilnice $\to$ Elemente de Itinerariu.
    \item  Many-to-Many: Locații $\leftrightarrow$ Etichete , relație rezolvată printr-o tabelă de joncțiune (LocationTag).
    \item  Preferințele utilizatorilor: Entitatea UserTagPreferences, implementată cu chei primare compuse (composite keys).
\end{itemize}

Configurația API-ului și a modelului fluent al EF Core gestionează aceste relații într-o manieră elegantă, asigurând managementul automat al cheilor externe și aplicarea riguroasă a politicilor de ștergere în cascadă.

Sistemul de migrări asigură controlul versiunilor pentru schema bazei de date. Această abordare este superioară utilizării scripturilor SQL manuale deoarece:

\begin{itemize}
    \item  migrările sunt integrate direct în limbaj și sigure din punct de vedere al tipurilor;
    \item  sistemul menține o pistă de audit completă a evoluției și mutațiilor schemei;
    \item  procedura de anulare a modificărilor este directă, predictibilă și complet reversibilă.
\end{itemize}

\subsection{Platfoma de Baze de Date: SQL Server}

SQL Server este alegerea ideală pentru aplicațiile ASP.NET datorită integrării native cu Visual Studio și Entity Framework, ceea ce simplifică procesul de dezvoltare. De asemenea garantează integritatea datelor, susține interogări analitice complexe, oferă și un mediu local extrem de eficient prin SQL Server LocalDB; care permite testarea și depanarea rapidă, direct pe stația de lucru.

\begin{figure}[htbp]
    \centering
    \includegraphics[width=0.7\linewidth]{Licenta.drawio.png}
    \caption{Figure}
    \label{fig:architecture_diagram}
\end{figure}

\section{Autentificare și Autorizare: ASP.NET Identity}

\subsection{Cadrul de gestionare a identității}
ASP.NET Core Identity este un sistem predefinit pentru gestionarea utilizatorilor. Adoptarea acestuia este optimă din următoarele considerente:

\begin{itemize}
    \item  hashing-ul și stocarea securizată a parolelor (folosind algoritmi de cripta adaptivi precum PBKDF2);
    \item  capacități de autentificare cu doi pași, în acest caz cu confirmarea prin e-mail;
    \item  controlul accesului bazat pe rol, activat prin intermediul AddRoles<IdentityRole>().
\end{itemize}

Integrare cu EF Core: Identity utilizează Entity Framework pentru persistența datelor, eliminând necesitatea unei biblioteci de autentificare separate și asigurând o integrare cu modelul de date al aplicației prin `IdentityDbContext<ApplicationUser>`.\href{https://arxiv.org/abs/2009.02114}{}\cite{identity_ref}

\section{Tehnologii Frontend}
\subsection{Framework-ul Bootstrap}
Bootstrap, inclus în structura proiectului sub directorul wwwroot/lib/bootstrap/, este un framework CSS esențial care furnizează un design adaptiv și asigură consistența vizuală a platformei, garantând o interfață uniformă pe toate paginile prin intermediul componentelor sale predefinite. În plus, sprijină accesibilitatea prin componentele sale implementate în conformitate cu directivele WCAG (Web Content Accessibility Guidelines). Totodată, accelerează dezvoltarea prin minimizarea codului CSS personalizat.

\subsection{jQuery și mecanismele de validare}
Biblioteca jQuery împreună cu extensiile sale dedicate validării oferă următoarele avantaje tehnice:
\begin{itemize}
    \item  validare non-intruzivă: regulile de validare sunt specificate în atributele data ale HTML, mai degrabă decât în JavaScript inline, păstrând preocupările separate;
    \item  prevenirea erorilor pe partea de client: înainte ca datele să ajungă la server, validarea pe partea clientului oferă feedback imediat utilizatorului, reducând încărcarea serverului și îmbunătățind experiența;
    \item  dacă JavaScript este dezactivat, validarea pe partea serverului funcționează în continuare, asigurând integritatea datelor.
\end{itemize}

Deși framework-urile moderne (React, Vue, Angular) ar putea înlocui jQuery, configurația actuală oferă o soluție ușoară, dovedită, care nu necesită complexitate suplimentară în fluxul de construire pentru o aplicație ASP.NET Core monolitică.

\section{Caracteristici geospațiale: Mapbox}
Mapbox este configurat ca furnizor principal de hărți\href{https://link.springer.com/chapter/10.1007/978-3-031-19594-5_2}{}\cite{mapbox_ref}. Pentru acest scenariu, Mapbox prezintă avantaje față de alte alternative deoarece:

\begin{itemize}
    \item  Eficiența costurilor: modelul de tarifare al platformei Mapbox se bazează pe utilizarea reală, ceea ce îl face rentabil pentru aplicațiile cu încărcături variabile de utilizatori. Vederile asociate permit randarea unor hărți interactive capabile să vizualizeze:
        \begin{itemize}
            \item  puncte de interes (POI) distribuite de-a lungul itinerariului;
            \item  rute și traiectorii topologice calculate între locații;
            \item  evidențierea țărilor vizitate.
        \end{itemize}
    \item  Prelucrarea datelor de geolocalizare: locațiile sunt persistente utilizând coordonate spațiale precise (latitudine și longitudine), capacitate care permite motorului Mapbox să efectueze operațiuni computaționale precum: calcularea precisă a distanțelor, optimizarea rutelor, reprezentarea vizuală a gradului de complexitate a topografiei.

\end{itemize}

\section{Servicii de e-mail: MailKit}
În contextul acestei aplicații, MailKit furnizează următoarele capabilități tehnice:

\begin{itemize}
    \item  suport extins pentru protocolul SMTP: integrarea directă cu infrastructura SMTP a platformei Gmail facilitează transmiterea automată a mesajelor pentru diferite necesități;
    \item  suportă standarde moderne de autentificare și criptare a canalelor de comunicație, garantând protecție;
    \item  serviciul EmailSender implementează interfața IEmailSender, permițând migrarea viitoare către serviciile cloud fără modificări de cod.
\end{itemize}

